\documentclass[11pt,letterpaper]{article}

% ── Packages ──────────────────────────────────────────────
\usepackage[margin=1in]{geometry}
\usepackage{amsmath,amssymb,amsthm}
\usepackage{mathtools}
\usepackage{booktabs}
\usepackage{enumitem}
\usepackage{hyperref}
\usepackage[dvipsnames]{xcolor}
\usepackage{titlesec}
\usepackage{microtype}
\usepackage{parskip}

% ── Formatting ────────────────────────────────────────────
\hypersetup{colorlinks=true, linkcolor=NavyBlue, citecolor=NavyBlue, urlcolor=NavyBlue}
\titleformat{\section}{\large\bfseries}{\thesection}{1em}{}
\titleformat{\subsection}{\normalsize\bfseries}{\thesubsection}{1em}{}
\titleformat{\subsubsection}{\normalsize\itshape}{\thesubsubsection}{1em}{}

% ── Custom environments ──────────────────────────────────
\newcommand{\caveat}[1]{\par\smallskip\noindent\textcolor{gray}{\textit{Caveat:} #1}\smallskip}
\newcommand{\pp}{\,\text{pp}}

% ── Math shortcuts ────────────────────────────────────────
\newcommand{\bfp}{\mathbf{p}}
\newcommand{\bfz}{\mathbf{z}}
\newcommand{\bfy}{\mathbf{y}}
\newcommand{\bfd}{\mathbf{d}}
\newcommand{\bfb}{\mathbf{b}}
\newcommand{\cS}{\mathcal{S}}
\newcommand{\cM}{\mathcal{M}}
\newcommand{\R}{\mathbb{R}}

\begin{document}

% ── Title ─────────────────────────────────────────────────
\begin{center}
{\LARGE\bfseries Cross-Patient Speech Decoding\\[4pt] from Intra-Operative $\mu$ECoG}\\[12pt]
{\large Theoretical Problem Framing}\\[8pt]
{\normalsize Ben Tang \quad\textbar\quad Greg Cogan Lab, Duke University}\\[4pt]
{\normalsize April 2026 \quad\textbar\quad Version 2}
\end{center}

\medskip\noindent\rule{\textwidth}{0.4pt}\medskip

\noindent This document formalizes the cross-patient speech decoding problem, organizes evidence into explicit epistemic categories, and identifies the structural bottlenecks that any solution must address. Claims are categorized as \textbf{externally established} (published literature), \textbf{internally observed} (our experiments), or \textbf{working hypotheses} (plausible but unverified). Architecture follows as a separate document.

% ══════════════════════════════════════════════════════════
\section{The Physical Problem}
% ══════════════════════════════════════════════════════════

Each patient's micro-electrocorticography ($\mu$ECoG) array is a \textbf{sparse spatial sample} of neural activity over the cortical surface during speech production. The array is placed subdurally over left ventral sensorimotor cortex (vSMC) during surgery, and its position in standardized brain space (MNI-152) is determined post-hoc via electrode reconstruction (BioImage Suite for DBS patients, Brainlab for tumor patients).

\subsection{The cortical surface, not free 3D space}

The cortex is a folded two-dimensional sheet, not a volume. Let $\cS$ denote the cortical surface embedded in $\R^3$. The high-gamma analytic amplitude (HGA, 70--150\,Hz envelope) is a function on this surface:
\begin{equation}
a: \cS \times \R \to \R
\end{equation}

For patient $i$ with $N_i$ electrodes at cortical surface positions $\{\bfp_j^{(i)}\}_{j=1}^{N_i} \subset \cS$, we observe:
\begin{equation}\label{eq:observation}
d_j^{(i)}(t) = a\bigl(\bfp_j^{(i)},\; t\bigr) + \epsilon_j^{(i)}(t)
\end{equation}
where $\epsilon$ captures measurement noise, electrode impedance variation, and sub-electrode-scale neural activity not resolved by the ${\sim}2$\,mm contact spacing.

In practice, we approximate cortical surface positions with 3D MNI coordinates: $\bfp_j^{(i)} \approx (x_j^{(i)}, y_j^{(i)}, z_j^{(i)})_{\text{MNI}}$. This approximation has two consequences:

\begin{enumerate}[leftmargin=2em]
\item \textbf{Euclidean distance in MNI space does not equal geodesic distance on the cortical surface.} Two electrodes 3\,mm apart in MNI may sit on opposite banks of a sulcus, making them functionally distant. For $\mu$ECoG arrays on exposed cortex, the local patch under the array is approximately planar, so within-array Euclidean distances are reasonable. Cross-patient comparisons at larger scales are more affected by folding geometry.

\item \textbf{Surface-based coordinates} (e.g., projection to \texttt{fsaverage}) would better respect cortical geometry than volumetric MNI coordinates. Whether surface reconstructions exist for our patients depends on the BioImage Suite\,/\,Brainlab outputs available.
\end{enumerate}

\textbf{The decoding problem.} Given observations $\{d_j^{(i)}(t)\}$ at approximately known surface positions, infer the phoneme sequence $\bfy = (y_1, y_2, y_3)$ that produced the spatiotemporal pattern.

\textbf{The cross-patient problem.} Train on patients $\{1, \ldots, K\}$ and decode for a held-out patient $K{+}1$ whose electrodes sample the cortical surface at \emph{different positions}. The observations are not aligned: electrode $j$ on patient $i$ and electrode $j$ on patient $i'$ correspond to different cortical locations.

\subsection{What makes this hard}

\begin{table}[h]
\centering\small
\begin{tabular}{@{}lll@{}}
\toprule
\textbf{Factor} & \textbf{Quantity} & \textbf{Implication} \\
\midrule
Electrodes per patient & 63--201 & Sparse sampling of the cortical field \\
Electrode spacing & ${\sim}2$\,mm & Resolves individual gyral columns \\
Array offset across patients & 15--25\,mm in MNI & Different cortical regions sampled \\
Array rotation & Variable & Channel indices meaningless across patients \\
Correspondence uncertainty & Several mm & Coordinate-based alignment is approximate \\
Cortical folding variability & Patient-specific & Sulcal geometry differs across individuals \\
Tumor mass effect (subset) & Variable & Macroscopic warping in 6 of 12 patients \\
Trials per patient & 46--178 & Small-data regime for supervised learning \\
Total utterance per patient & ${\sim}1$\,min & Insufficient for single-patient SSL \\
Phoneme vocabulary & 9 classes $\times$ 3 positions & 27-way classification at each trial \\
\bottomrule
\end{tabular}
\end{table}

% ══════════════════════════════════════════════════════════
\section{The Shared Dynamics Hypothesis}
% ══════════════════════════════════════════════════════════

Cross-patient decoding is only possible if the cortical field $a(\bfp, t)$ shares structure across patients. We decompose:
\begin{equation}\label{eq:decomposition}
a^{(i)}(\bfp, t) = f\bigl(T_i(\bfp),\; t;\; \bfy\bigr) + r^{(i)}(\bfp, t)
\end{equation}
where:
\begin{itemize}[leftmargin=2em]
\item $f(\cdot, t; \bfy)$ is the \textbf{shared neural field}---the spatiotemporal pattern that phoneme sequence $\bfy$ produces in canonical cortical space.
\item $T_i: \cS_i \to \cS_{\text{canonical}}$ is a \textbf{patient-specific spatial transform} mapping each patient's cortical surface to a canonical reference. MNI normalization is our estimate of $T_i$, but it is imperfect.
\item $r^{(i)}(\bfp, t)$ is a \textbf{patient-specific residual} capturing individual neural organization, signal characteristics, and functional plasticity.
\end{itemize}
The viability of cross-patient decoding depends on the ratio $\|f\| / \|r^{(i)}\|$ at the spatial resolution accessible to our measurements and alignment procedure.

\subsection{Evidence for shared dynamics}

We organize evidence into three explicit categories.

\subsubsection{Externally established (published literature)}

\textbf{Per-patient input layers + shared temporal backbone is the optimal transfer architecture.} Singh et al.\ (2025) demonstrate group training with $N{=}25$ sEEG patients: per-patient Conv1D + shared BiLSTM ($2{\times}64$), frozen shared backbone + fine-tuned per-patient layers yields PER 0.49 vs.\ 0.57 single-subject ($p{<}0.001$). Willett et al.\ (2023) use per-day affine input layer ($256{\to}256$) + shared 5-layer GRU (512 hidden) for day-to-day BCI transfer.

\caveat{Singh's data regime is ${\sim}60$\,min/patient (vs.\ our ${\sim}1$\,min); Conv1D hyperparameters are unreported. Willett is $N{=}1$ with ${\sim}40$\,min/day. The architectural pattern transfers; the hyperparameters do not.}

\textbf{Articulatory somatotopy in vSMC is reproducible within individuals.} Metzger et al.\ (2023) show somatotopic organization of speech articulators in one patient with chronic ECoG (253 channels): labial, front tongue, back tongue, and vocalic representations are spatially organized ($p{<}0.0001$), preserved even after 18 years of anarthria, suggesting the organization is intrinsic to motor cortex circuitry.

\caveat{$N{=}1$ chronic implant. Cross-patient spatial consistency is inferred from anatomical uniformity but not directly measured in any paper in our corpus.}

\textbf{Articulatory features are a robust encoding basis in vSMC.} Wu et al.\ (2025) reconstruct speaker-invariant articulatory features from HD-ECoG over vSMC (PCC 0.75--0.80, $N{=}3$), including one intra-operative patient (PCC\,=\,0.76). The claim that articulatory features are ``more robustly encoded than acoustic features'' is attributed by Wu to Chartier (2018), Conant (2018), and Anumanchipalli (2019).

\caveat{Wu's reconstruction is word-level scalar loadings, not continuous trajectories. The robustness claim is borrowed from cited literature.}

\textbf{MNI coordinates enable cross-patient ECoG transfer.} Chen et al.\ (2025, SwinTW) use MNI coordinates and ROI embeddings as positional information across 43 ECoG + 9 sEEG subjects, reporting mean PCC\,=\,0.765 on unseen participants in cross-validation.

\caveat{Standard ECoG has ${\sim}10$\,mm spacing, comparable to inter-subject correspondence uncertainty. Our $\mu$ECoG spacing (${\sim}2$\,mm) is finer, meaning coordinates may encode global array position but not resolve within-array relationships.}

\textbf{Motor cortex neural manifold structure is preserved across individuals.} Safaie et al.\ (2023, \textit{Nature Neuroscience}) show that motor cortex population dynamics during reaching are geometrically similar across different monkeys.

\caveat{Limb motor cortex in non-human primates, not speech motor cortex in humans. Extension to speech is plausible but not directly tested.}

\textbf{Neural manifold dynamics are stable despite neural turnover.} Gallego et al.\ (2020, \textit{Nature Neuroscience}) establish that low-dimensional manifold dynamics are preserved across long time periods within the same animals, even as recorded neurons change. This is evidence for \emph{temporal stability} of latent dynamics, distinct from the cross-individual claim above.

\caveat{Within-individual stability, not cross-individual similarity.}

\textbf{Consistent articulatory somatotopy across participants.} Bouchard et al.\ (2013, \textit{Nature}) demonstrated reproducible articulatory maps in vSMC across 3 participants. Anumanchipalli et al.\ (2019, \textit{Nature}) reconstructed continuous articulatory trajectories from motor cortex.

\caveat{Neither paper is summarized in our corpus. Bouchard's $N{=}3$ provides limited evidence for population-level consistency.}

\subsubsection{Internally observed (our experiments)}

\textbf{Cross-patient supervised transfer works at this data scale.} Our LOPO pipeline (per-patient SpatialConv + shared BiGRU backbone, 9 source patients) achieves PER\,=\,0.762 on held-out S14, compared to per-patient PER\,=\,0.737 (grouped-by-token CV). This is direct evidence that decodable shared structure exists in~$f$.

\textbf{Articulatory bottleneck head improves LOPO decoding.} Projecting to a 15-dimensional phonological feature space via a fixed linguistic matrix outperforms a flat CE head for cross-patient decoding, consistent with $f$ being organized along articulatory feature dimensions.

\textbf{55 experiments on S14 converge to PER 0.750--0.780.} Across supervised, SSL, domain adaptation, self-training, and knowledge distillation paradigms. See Section~\ref{sec:wall} for interpretation.

\textbf{SSL features carry partial structure.} JEPA features + LP-FT achieve PER\,=\,0.797 (vs.\ 0.854 basic evaluation). Supervised features with the same recipe achieve 0.737.

\subsubsection{Working hypotheses (plausible but unverified)}

The ${\sim}0.76$ PER convergence on S14 reflects (in part) a spatial alignment bottleneck. This has not been isolated from target-data scarcity, measurement noise, or S14-specific effects.

MNI coordinate-based spatial alignment will improve decoding. The practical gain depends on registration accuracy, model architecture, and whether spatial alignment is binding.

Cross-task data pooling (Duraivel 2025 pseudoword repetition, same 9 phonemes, 3 overlapping patients) is a high-value, low-risk path to more data.


% ══════════════════════════════════════════════════════════
\section{The Neural Manifold Connection}
% ══════════════════════════════════════════════════════════

\subsection{Two views of the same phenomenon}

The cortical field model (Section~2) and the neural manifold hypothesis describe the same system at different levels:
\begin{itemize}[leftmargin=2em]
\item \textbf{The field view} (spatial): Activity $a(\bfp, t)$ is a function on the cortical surface. Each electrode samples this field at a known position.
\item \textbf{The manifold view} (population): The vector of all electrode activities $\bfd(t) = [d_1(t), \ldots, d_N(t)]$ lies on a low-dimensional manifold $\cM \subset \R^N$. Motor behavior corresponds to trajectories on $\cM$.
\end{itemize}

The bridge is a \textbf{latent state variable}. Let $\bfz(t) \in \R^k$ be the low-dimensional articulatory state---the configuration of the vocal tract at time $t$. Each electrode's activity is determined by \emph{where it sits} on the cortical surface and \emph{what the articulatory state is}:
\begin{equation}\label{eq:obs-function}
d_j^{(i)}(t) = g\bigl(\bfp_j^{(i)},\; \bfz(t)\bigr) + \epsilon_j^{(i)}(t)
\end{equation}
where $g: \cS \times \R^k \to \R$ is a \textbf{shared observation function} mapping (cortical position, articulatory state) to neural activity. This is the somatotopic map made explicit: electrodes over lip motor cortex respond to lip-state variables; electrodes over tongue motor cortex respond to tongue-state variables.

The dynamics on the manifold are determined by the phoneme sequence:
\begin{equation}\label{eq:dynamics}
\dot{\bfz}(t) = h\bigl(\bfz(t);\; \bfy\bigr)
\end{equation}

\subsection{Why this formulation matters}

The observation function $g(\bfp, \bfz)$ is the central learning target. It says: \textbf{the same function generates neural activity for all patients.} What differs is the set of positions $\{\bfp_j^{(i)}\}$ at which each patient samples $g$.

Without knowing electrode positions, the model must learn a separate observation function per patient---this is what per-patient SpatialConv does. With positions (MNI coordinates as an approximation of $\bfp$), the observation function~$g$ can be parameterized as a shared function of brain coordinates, and cross-patient alignment is built into the architecture.

\subsection{The articulatory manifold has known structure}

For speech, the manifold $\cM$ has physical meaning corresponding to vocal tract configuration:

\begin{table}[h]
\centering\small
\begin{tabular}{@{}lcl@{}}
\toprule
\textbf{Articulator} & \textbf{Approx.\ DoF} & \textbf{Controls} \\
\midrule
Tongue (tip, body, dorsum) & ${\sim}6$ & Place of articulation, vowel quality \\
Jaw & ${\sim}2$ & Aperture, protrusion \\
Lips & ${\sim}2$ & Rounding, closure \\
Larynx & ${\sim}2$ & Voicing, pitch \\
Velum & ${\sim}1$ & Nasal/oral \\
\midrule
\textbf{Total} & ${\sim}$\textbf{13--15} & \\
\bottomrule
\end{tabular}
\end{table}

For our 9 phonemes, the effective dimensionality is lower:

\begin{table}[h]
\centering\small
\begin{tabular}{@{}llc@{}}
\toprule
\textbf{Contrast} & \textbf{Phonemes distinguished} & \textbf{Approx.\ dim.} \\
\midrule
Place of articulation & labial (b,p,v) vs.\ velar (g,k) & ${\sim}2$ \\
Manner & stop vs.\ fricative vs.\ vowel & ${\sim}2$ \\
Voicing & voiced (b,g,v) vs.\ voiceless (k,p) & ${\sim}1$ \\
Vowel height & low (a,\ae) vs.\ high (i,u) & ${\sim}1$ \\
Vowel backness & front (\ae,i) vs.\ back (a,u) & ${\sim}1$ \\
\bottomrule
\end{tabular}
\end{table}

\noindent\textbf{Effective manifold dimension for our task: ${\sim}$5--8.} This is far below the electrode count (63--201), consistent with massive oversampling of a low-dimensional manifold. The decoding problem reduces to: given sparse samples of $g(\bfp, \bfz)$ at known positions~$\bfp$, infer the latent trajectory $\bfz(t)$ and classify the phoneme sequence.

\subsection{Cross-patient transfer as manifold alignment}

Each patient's electrode array gives a different \textbf{projection} of the shared manifold:
\begin{equation}\label{eq:projection}
\bfd^{(i)}(t) \approx W^{(i)} \bfz(t) + \boldsymbol{\epsilon}^{(i)}
\end{equation}
where $W^{(i)} \in \R^{N_i \times k}$ is determined by which cortical positions the electrodes sample.

Existing alignment methods (CCA, Procrustes, CORAL) attempt to align these projections post-hoc. All failed in our experiments. MNI coordinates offer an alternative: \textbf{aligning the projections a priori} by providing the geometric relationship between electrode positions.

The key question is whether inter-subject correspondence uncertainty (several mm) is small enough relative to the somatotopic map scale (${\sim}15$--$20$\,mm) that MNI coordinates provide useful alignment. Chen et al.\ (2025) suggest yes for standard ECoG. Whether this holds for $\mu$ECoG at ${\sim}2$\,mm spacing is an open empirical question.


% ══════════════════════════════════════════════════════════
\section{Decomposing Cross-Patient Variation}
% ══════════════════════════════════════════════════════════

\subsection{The spatial transform $T_i$}

The primary source of cross-patient variation is \textbf{geometric}: different electrode positions sample different regions of the shared field. The transform $T_i$ decomposes:
\begin{equation}\label{eq:transform}
T_i(\bfp) = A_i \cdot \bfp + \bfb_i + \delta_i(\bfp)
\end{equation}
where:
\begin{itemize}[leftmargin=2em]
\item $A_i \in \R^{3 \times 3}$: rigid or affine transform (rotation, scaling). MNI normalization estimates and corrects this.
\item $\bfb_i \in \R^3$: translation (array placement offset). Our arrays span 15--25\,mm in MNI.
\item $\delta_i(\bfp)$: nonlinear residual---sulcal geometry, local cortical folding, and for tumor patients, \textbf{macroscopic warping from mass effect}. Tumor displacement is a geometric mapping error, not signal noise: it physically moves the somatotopic map relative to skull-based coordinates. This is the most severe form of $\delta_i$ in our cohort and violates the affine assumption of $A_i$.
\end{itemize}

\subsection{Inter-subject correspondence uncertainty}

The commonly cited ``${\sim}5$\,mm'' figure is not a stable physical constant. The actual uncertainty has multiple sources:

\begin{enumerate}[leftmargin=2em]
\item \textbf{Native-space electrode localization}: Validated pipelines can achieve sub-millimeter accuracy---this step is relatively precise.
\item \textbf{Brain shift correction}: Intra-operative brain shift (craniotomy, CSF drainage, gravity) displaces the cortex relative to pre-operative imaging. Correction quality varies.
\item \textbf{Warping to common atlas space}: The nonlinear registration from patient anatomy to MNI-152 is the dominant source of inter-subject uncertainty. Quality depends on algorithm, template, and---for tumor patients---whether the lesion distorts the warp.
\end{enumerate}

The net result is \textbf{several millimeters of inter-subject correspondence uncertainty} for healthy anatomy, potentially more near tumors:
\begin{itemize}[leftmargin=2em]
\item At the scale of articulatory somatotopy (${\sim}15$--$20$\,mm zones): correspondence is reliable.
\item At the scale of electrode spacing (${\sim}2$\,mm): correspondence is unreliable.
\item In between (${\sim}5$--$10$\,mm): partially informative, varies by patient and region.
\end{itemize}

\subsection{Resolution mismatch and dual spatial encoding}

$\mu$ECoG electrode spacing (${\sim}2$\,mm) is finer than inter-subject correspondence uncertainty. This implies any coordinate-aware architecture should use \textbf{dual spatial encoding}:

\begin{table}[h]
\centering\small
\begin{tabular}{@{}llll@{}}
\toprule
\textbf{Scale} & \textbf{Resolution} & \textbf{Best encoded by} & \textbf{What it captures} \\
\midrule
Cross-patient & ${\sim}10$--$15$\,mm & MNI coordinates & Articulatory somatotopy \\
Intra-patient global & ${\sim}5$--$10$\,mm & MNI (approximate) & Array position on cortex \\
Intra-patient local & ${\sim}2$\,mm & Grid topology & Fine-grained spatial patterns \\
\bottomrule
\end{tabular}
\end{table}

\subsection{The patient-specific residual $r^{(i)}$}

After accounting for shared dynamics and spatial transform, the residual captures:
\begin{enumerate}[leftmargin=2em]
\item \textbf{Fine-grained neural organization}: Within-articulator spatial patterns that vary across individuals.
\item \textbf{Signal characteristics}: Electrode impedance, contact quality, and SNR. Z-score normalization addresses amplitude scaling but not higher-order distributional differences.
\item \textbf{Functional plasticity}: Individual speech development, accent, and articulatory habits.
\item \textbf{Pathology-specific effects}: Parkinson's motor cortex excitability changes; tumor-adjacent cortical reorganization (distinct from the geometric displacement captured in $\delta_i$).
\end{enumerate}


% ══════════════════════════════════════════════════════════
\section{What Our Experiments Reveal}
% ══════════════════════════════════════════════════════════

\subsection{The evaluation recipe dominates at this data scale}

Across experiments on S14, training improvements contributed ${\sim}7\pp$ and evaluation improvements contributed ${\sim}6.5\pp$:

\begin{table}[h]
\centering\small
\begin{tabular}{@{}llr@{}}
\toprule
\textbf{Component} & \textbf{Category} & \textbf{PER improvement} \\
\midrule
Label smoothing 0.1 & Training & $-4.1\pp$ \\
Mixup $\alpha{=}0.2$ & Training & $-1.7\pp$ \\
Per-position heads + dropout & Training & $-1.2\pp$ \\
\midrule
Weighted $k$-NN ($k{=}10$, cosine) & Evaluation & $-2.1\pp$ \\
TTA (16 augmented copies) & Evaluation & $-1.4\pp$ \\
Articulatory + CE dual head & Evaluation & $-1.1\pp$ \\
Best-of per fold & Evaluation & $-1.8\pp$ \\
\bottomrule
\end{tabular}
\end{table}

\subsection{Multi-scale temporal structure is validated}

Multi-scale temporal encoding (stride 3+5+10) was the single best architectural improvement (PER $0.766 \to 0.750$):

\begin{table}[h]
\centering\small
\begin{tabular}{@{}lll@{}}
\toprule
\textbf{Scale} & \textbf{Temporal resolution} & \textbf{Articulatory correlate} \\
\midrule
Stride 3 (${\sim}67$\,Hz) & ${\sim}15$\,ms & Fast transitions (stop release, tongue flap) \\
Stride 5 (${\sim}40$\,Hz) & ${\sim}25$\,ms & Syllable dynamics (CV transition) \\
Stride 10 (${\sim}20$\,Hz) & ${\sim}50$\,ms & Phoneme-level patterns (steady-state) \\
\bottomrule
\end{tabular}
\end{table}

\subsection{SSL features carry partial structure}

All SSL methods produce near-chance PER (0.85--0.91) under basic evaluation. JEPA features + LP-FT achieve PER\,=\,0.797, and weighted $k$-NN improves them by $2.9\pp$. SSL features capture some structure but their feature geometry is far weaker than supervised features. TTA hurts JEPA features ($-0.2\pp$), confirming they are not augmentation-invariant.

\subsection{The ${\sim}$0.76 PER convergence on S14}\label{sec:wall}

55 experiments across all paradigms converge to PER 0.750--0.780 on S14. \textbf{Attributing this to any single bottleneck would be premature}:

\begin{enumerate}[leftmargin=2em]
\item \textbf{S14 is no longer a held-out test patient.} With 55+ architectures evaluated against S14, it functions as a validation set. The convergence may partly reflect S14-specific characteristics rather than a universal ceiling.
\item \textbf{Measurement noise from small validation folds.} ${\sim}30$ samples per fold implies ${\pm}0.05$ PER noise. Differences ${<}2\pp$ are within noise.
\item \textbf{Target-data scarcity.} S14 has 153 trials (${\sim}120$ per training fold). Stage~2 adaptation contributes only ${\sim}1.6\pp$.
\item \textbf{Spatial alignment may be one contributing factor.} Per-patient spatial encoders cannot share spatial knowledge. But this has not been isolated as \emph{the} bottleneck.
\end{enumerate}

Breaking through requires testing multiple hypotheses: \textbf{(a)}~population-level evaluation on all patients, \textbf{(b)}~cross-task data pooling, and \textbf{(c)}~MNI coordinate-based spatial alignment. These are complementary, not competing.


% ══════════════════════════════════════════════════════════
\section{The Case for MNI Coordinates}
% ══════════════════════════════════════════════════════════

\subsection{Currently unused information}

The current LOPO pipeline processes spatial information through \textbf{per-patient Conv2d} operating on grid topology:
\[
d^{(i)} \in \R^{H_i \times W_i \times T}
\xrightarrow{\;\text{Conv2d}_i\;}
\R^{d \times T}
\xrightarrow{\;\text{shared backbone}\;}
\R^{d' \times T'}
\]

Each patient's Conv2d is independent. The model has \textbf{zero information} about the spatial relationship between electrodes across patients. The shared backbone implicitly assumes its input features are homologous across the channel dimension---that feature dimension $k$ from different patients' SpatialConv layers carries corresponding information. This assumption is anatomically false: independent per-patient Conv2d layers produce feature spaces with no guaranteed alignment.

\subsection{What coordinates enable (in principle)}

The observation function formulation~\eqref{eq:obs-function} makes the value explicit: knowing $\bfp_j^{(i)}$ allows the model to learn $g$ as a shared function rather than learning independent mappings per patient. Concretely:

\begin{enumerate}[leftmargin=2em]
\item \textbf{Identifying which articulatory zone each electrode samples.} The somatotopic map operates at ${\sim}15$--$20$\,mm scale, above the inter-subject correspondence uncertainty.
\item \textbf{Shared spatial representations across patients.} Electrodes at similar MNI coordinates can share representations, pooling spatial information across the cohort.
\item \textbf{Contextualizing variable electrode counts.} Coordinates tell the model which cortical regions are covered and which are absent for each patient.
\end{enumerate}

\subsection{What coordinates do NOT guarantee}

The availability of MNI coordinates does not automatically improve decoding. Coordinates help only if:
\begin{itemize}[leftmargin=2em]
\item The correspondence uncertainty is small enough relative to the somatotopic scale (plausible for ${\sim}15$--$20$\,mm zones, uncertain within ${\sim}5$\,mm neighborhoods).
\item The model architecture can exploit the information (a design question).
\item Spatial alignment is actually a binding constraint on performance (unverified).
\item The coordinates are accurate enough for specific patients (tumor patients may have larger errors).
\end{itemize}

\noindent\textbf{MNI coordinates are a necessary but not sufficient ingredient for coordinate-based cross-patient transfer.} Their value must be demonstrated empirically.


% ══════════════════════════════════════════════════════════
\section{Available Data and Experimental Paths}
% ══════════════════════════════════════════════════════════

\subsection{Currently used}

\begin{table}[h]
\centering\small
\begin{tabular}{@{}lll@{}}
\toprule
\textbf{Data source} & \textbf{Volume} & \textbf{Usage} \\
\midrule
PS task trial-level HGA & 9 patients, ${\sim}1315$ trials & Supervised LOPO training \\
S14 trial-level HGA & 153 trials & Target patient evaluation \\
Grid topology (from TSV) & Per-patient & SpatialConv input structure \\
\bottomrule
\end{tabular}
\end{table}

\subsection{Available but unused}

\begin{table}[h]
\centering\small
\begin{tabular}{@{}lllc@{}}
\toprule
\textbf{Data source} & \textbf{Volume} & \textbf{Potential usage} & \textbf{Risk} \\
\midrule
Cross-task pooling (Duraivel) & 3 patients, 156--208 trials & Doubles labeled data & Low \\
MNI coordinates & All patients & Spatial alignment & Medium \\
Continuous recordings & ${\sim}168$\,min & Temporal SSL & Medium \\
Lexical task patients & 13 patients & Additional sources & Medium \\
Audio recordings & 4 variants/patient & Auxiliary supervision & Medium \\
\bottomrule
\end{tabular}
\end{table}

\subsection{Cross-task pooling deserves priority}

The pseudoword repetition data (Duraivel 2025) is the lowest-risk path to more training data: same 52 CVC/VCV tokens, same 9 phonemes, 3 overlapping patients, same preprocessing pipeline, same IRB. This should be pursued as a \textbf{central experimental branch}, not a side opportunity.

\subsection{The continuous data caveat}

The ${\sim}168$ minutes of continuous recordings are not ${\sim}168$ minutes of useful SSL signal. Intra-operative recordings are dominated by silence, operating room ambient noise, auditory perception, and anesthesia-related fluctuations. Utilizing this data for SSL requires aggressive curation (VAD, energy thresholding) to extract speech-relevant neural activity.


% ══════════════════════════════════════════════════════════
\section{Summary}
% ══════════════════════════════════════════════════════════

\subsection*{Externally established}
\begin{itemize}[leftmargin=2em]
\item Per-patient spatial layer + shared temporal backbone is optimal (Singh 2025, Willett 2023)
\item Articulatory somatotopy exists in vSMC and is intrinsic to circuitry (Metzger 2023, Bouchard 2013)
\item MNI/atlas coordinates enable cross-patient ECoG transfer (Chen 2025)
\item Motor cortex manifold dynamics are preserved across individuals (Safaie 2023) and stable over time (Gallego 2020)
\item Articulatory features are robustly encoded in vSMC (Wu 2025, citing Chartier 2018, Anumanchipalli 2019)
\end{itemize}

\subsection*{Internally observed}
\begin{itemize}[leftmargin=2em]
\item Cross-patient LOPO transfer works: PER 0.762 vs.\ per-patient 0.737 (S14, grouped CV)
\item Articulatory bottleneck head improves cross-patient decoding
\item Multi-scale temporal encoding (stride 3+5+10) is the best validated architectural improvement
\item Evaluation recipe contributes ${\sim}6.5\pp$, rivaling training improvements (${\sim}7\pp$)
\item All SSL methods fail at current data scale but features carry partial structure
\item 55 experiments on S14 converge to PER 0.750--0.780 across all paradigms
\item All domain adaptation methods (DANN, CORAL, CCA) provide zero benefit
\end{itemize}

\subsection*{Working hypotheses}
\begin{itemize}[leftmargin=2em]
\item The ${\sim}0.76$ convergence reflects (in part) a spatial alignment bottleneck---not yet isolated from data scarcity, noise, or S14-specific effects
\item MNI coordinate-based alignment will improve decoding---depends on registration accuracy, architecture, and whether alignment is binding
\item The neural field decomposes into shared dynamics + spatial transform + residual---productive but magnitudes unmeasured
\item Cross-task pooling is high-value, low-risk---supported by task compatibility but untested in pipeline
\item Continuous recordings provide SSL signal after curation---usable fraction unknown
\end{itemize}

\subsection*{The central question}

Our current results are consistent with the hypothesis that explicit spatial correspondence is one remaining bottleneck, but they do not yet isolate that bottleneck from target-data scarcity, fold noise, and coverage mismatch. The highest-confidence experimental paths are: \textbf{(a)}~population-level evaluation on all patients, \textbf{(b)}~cross-task data pooling, and \textbf{(c)}~MNI coordinate-based spatial alignment. These address complementary aspects of the problem and should be pursued in parallel where feasible.

\end{document}
